\section{Conclusão}

Este trabalho percorreu o caminho entre um estudo de aprendizado de máquina em lote
e um sistema de detecção de intrusão que opera sobre tráfego real. No MVP1, a
análise da base NSL-KDD confirmou o que a literatura relata: modelos baseados em
árvore atingem acurácia acima de 99,7\%, e 17 atributos bastam para isso, com a
\en{Decision Tree} oferecendo o melhor equilíbrio entre acurácia e latência de
predição. O mesmo estudo mostrou, porém, que essa acurácia não sobrevive a ataques
ausentes do treino: no \cod{KDDTest+}, o \en{recall} caiu para 63,22\%, com o erro
concentrado em falsos negativos silenciosos.

Essa lacuna orientou o MVP2. O sistema resultante combina um detector não
supervisionado (\en{Isolation Forest}), que não depende de ter visto o ataque, com
um refinador supervisionado (SVM linear), que aprende com os vereditos do
analista, e uma árvore \en{surrogate} que traduz cada alerta em uma regra legível.
As medições mostraram que os dois estágios são de fato complementares: o
\en{Isolation Forest} tem \en{recall} nulo contra força bruta real (FTP-Patator),
por uma razão estrutural --- ataques volumosos e homogêneos não são raros ---, e o
SVM recupera esses ataques quando os conhece.

A principal contribuição do trabalho está no ciclo de vida do modelo. Retreinar um
detector de anomalias sobre o tráfego recente o torna explorável por um atacante
paciente, e as métricas mais óbvias para vigiar o retreino falham justamente nesse
caso: a fidelidade da explicação não se move e a taxa de falso positivo até
melhora enquanto o detector fica cego. A resposta adotada --- \en{pool} curado com
quarentena, recência ponderada, despejo com procedência e um portão que mede
\en{recall} por cenário com piso absoluto --- reprova candidatos envenenados a
partir de 0,10\% do \en{pool} e limita o dano acumulado ao intervalo entre a
produção e o piso. O sistema foi entregue conteinerizado, sem privilégio de
\en{root}, com persistência em TimescaleDB e painéis em Grafana, e verificado de
ponta a ponta.

Como trabalhos futuros, destacam-se: (i) ampliar o \en{golden set} com capturas
próprias de testes de intrusão para todos os cenários, reduzindo a variância do
cenário canário; (ii) implementar atributos de horizonte longo, agregados por par
origem/destino ao longo de dias, para detectar beacons de período longo; (iii)
medir o Estágio 2 com vereditos reais e validação temporal; (iv) concluir a
interface de veredito do analista e a edição de parâmetros pelo Grafana; e (v)
operar um segundo canal de detecção por assinatura (Suricata ou Zeek) em paralelo,
para comparação e correlação de alertas.

\subsection{Cronograma}

\begin{enumerate}
	\item \textbf{Definição do tema:} estudos introdutórios sobre IDS e aplicação de ML na detecção de ameaças.
	\item \textbf{Revisão bibliográfica:} levantamento de trabalhos sobre detecção de intrusões, pré-processamento e algoritmos de ML.
	\item \textbf{MVP1:} estudo em lote sobre a NSL-KDD, comparação de algoritmos e redução de dimensionalidade.
	\item \textbf{Apresentação parcial:} apresentação do MVP1 e da fundamentação teórica.
	\item \textbf{MVP2:} pipeline em dois estágios, ciclo de vida, persistência, contêineres e painéis.
	\item \textbf{Entrega final:} relatório completo e apresentação do sistema.
\end{enumerate}

\begin{strip}
\centering
\footnotesize
\renewcommand{\arraystretch}{0.9}
\setlength{\tabcolsep}{2pt}
\begin{tabular}{|p{2.8cm}|*{12}{c|}}
\hline
\multirow{2}{*}{\textbf{Item}} & \multicolumn{12}{c|}{\textbf{2026}} \\ \cline{2-13}
& Jan & Fev & Mar & Abr & Mai & Jun & Jul & Ago & Set & Out & Nov & Dez \\
\hline
Definição do tema       & X & X &   &   &   &   &   &   &   &   &   &   \\ \hline
Revisão bibliográfica   &   & X & X & X &   &   &   &   &   &   &   &   \\ \hline
Desenvolvimento MVP1    &   &   &   & X & X &   &   &   &   &   &   &   \\ \hline
Apresentação parcial    &   &   &   &   &   & X &   &   &   &   &   &   \\ \hline
Desenvolvimento MVP2    &   &   &   &   &   &   & X & X & X & X & X &   \\ \hline
Entrega do TCC          &   &   &   &   &   &   &   &   &   &   &   & X \\ \hline
\end{tabular}
\captionof{table}{Cronograma do projeto.}
\end{strip}
